@c -*-texinfo-*-

@c  Copyright (c) 1996 Blake McBride
@c  All rights reserved.
@c
@c  Redistribution and use in source and binary forms, with or without
@c  modification, are permitted provided that the following conditions are
@c  met:
@c
@c  1. Redistributions of source code must retain the above copyright
@c  notice, this list of conditions and the following disclaimer.
@c
@c  2. Redistributions in binary form must reproduce the above copyright
@c  notice, this list of conditions and the following disclaimer in the
@c  documentation and/or other materials provided with the distribution.
@c
@c  THIS SOFTWARE IS PROVIDED BY THE COPYRIGHT HOLDERS AND CONTRIBUTORS
@c  "AS IS" AND ANY EXPRESS OR IMPLIED WARRANTIES, INCLUDING, BUT NOT
@c  LIMITED TO, THE IMPLIED WARRANTIES OF MERCHANTABILITY AND FITNESS FOR
@c  A PARTICULAR PURPOSE ARE DISCLAIMED. IN NO EVENT SHALL THE COPYRIGHT
@c  HOLDER OR CONTRIBUTORS BE LIABLE FOR ANY DIRECT, INDIRECT, INCIDENTAL,
@c  SPECIAL, EXEMPLARY, OR CONSEQUENTIAL DAMAGES (INCLUDING, BUT NOT
@c  LIMITED TO, PROCUREMENT OF SUBSTITUTE GOODS OR SERVICES; LOSS OF USE,
@c  DATA, OR PROFITS; OR BUSINESS INTERRUPTION) HOWEVER CAUSED AND ON ANY
@c  THEORY OF LIABILITY, WHETHER IN CONTRACT, STRICT LIABILITY, OR TORT
@c  (INCLUDING NEGLIGENCE OR OTHERWISE) ARISING IN ANY WAY OUT OF THE USE
@c  OF THIS SOFTWARE, EVEN IF ADVISED OF THE POSSIBILITY OF SUCH DAMAGE.



@section Tree Views
The @code{TreeView} class (a subclass of @code{Control}) implements a
hierarchical tree control for displaying and selecting items arranged in
a parent-child structure.  Individual nodes in the tree are represented
by instances of the @code{TreeViewItem} class.

@subsection TreeView Methods



@deffn {AddTVItem} gAddTVItem::TreeView
@findex gAddTVItem
@sp 2
@example
@group
tvi = gAddTVItem(tv, parent, txt);

object  tv;       /*  TreeView object          */
object  parent;   /*  parent TreeViewItem or NULL  */
char    *txt;     /*  item display text         */
object  tvi;      /*  new TreeViewItem object   */
@end group
@end example
Creates and adds a new tree item with the given text.  If @code{parent}
is @code{NULL}, the item is added at the root level.  Otherwise it is
added as a child of @code{parent}.  Returns the newly created
@code{TreeViewItem}.
@sp 1
See also:  @code{AddTVItemWithImage}
@end deffn



@deffn {AddTVItemWithImage} gAddTVItemWithImage::TreeView
@findex gAddTVItemWithImage
@sp 2
@example
@group
tvi = gAddTVItemWithImage(tv, parent, image, txt);

object  tv;       /*  TreeView object              */
object  parent;   /*  parent TreeViewItem or NULL   */
int     image;    /*  image index in image list     */
char    *txt;     /*  item display text             */
object  tvi;      /*  new TreeViewItem object       */
@end group
@end example
Creates and adds a new tree item with the given text and a custom image
from the associated image list.  Similar to @code{gAddTVItem} but allows
specifying an image index.
@sp 1
See also:  @code{AddTVItem}, @code{SetImageList}, @code{NewImageList}
@end deffn



@deffn {Alphabetize} gAlphabetize::TreeView
@findex gAlphabetize::TreeView
@sp 2
@example
@group
r = gAlphabetize(tv);

object  tv;   /*  TreeView object  */
object  r;    /*  TreeView object  */
@end group
@end example
Enables alphabetical sorting of tree items.
@sp 1
See also:  @code{AlphabetizeOff}, @code{Resort}
@end deffn



@deffn {AlphabetizeOff} gAlphabetizeOff::TreeView
@findex gAlphabetizeOff::TreeView
@sp 2
@example
@group
r = gAlphabetizeOff(tv);

object  tv;   /*  TreeView object  */
object  r;    /*  TreeView object  */
@end group
@end example
Disables alphabetical sorting.  Items will sort by their original
insertion order (or by @code{order}/@code{suborder} properties).
@sp 1
See also:  @code{Alphabetize}, @code{Resort}
@end deffn



@deffn {CheckFunction} gCheckFunction::TreeView
@findex gCheckFunction::TreeView
@sp 2
@example
@group
r = gCheckFunction(tv, fun);

object  tv;    /*  TreeView object              */
int     (*fun)();  /*  validation function       */
object  r;     /*  TreeView object              */
@end group
@end example
Sets a validation function that is called by @code{gCheckValue} to
validate the tree view's state before the dialog completes.
@sp 1
See also:  @code{CheckValue}
@end deffn



@deffn {CheckValue} gCheckValue::TreeView
@findex gCheckValue::TreeView
@sp 2
@example
@group
err = gCheckValue(tv);

object  tv;    /*  TreeView object           */
int     err;   /*  0=valid, 1=error          */
@end group
@end example
Validates the tree view's current state using the function set by
@code{gCheckFunction}.  Returns 0 if valid, 1 if validation fails.
@sp 1
See also:  @code{CheckFunction}
@end deffn



@deffn {Deselect} gDeselect::TreeView
@findex gDeselect::TreeView
@sp 2
@example
@group
r = gDeselect(tv);

object  tv;   /*  TreeView object  */
object  r;    /*  TreeView object  */
@end group
@end example
Deselects the current selection in the tree view.
@end deffn



@deffn {First} gFirst::TreeView
@findex gFirst::TreeView
@sp 2
@example
@group
tvi = gFirst(tv);

object  tv;    /*  TreeView object             */
object  tvi;   /*  first TreeViewItem or NULL   */
@end group
@end example
Returns the root (first) @code{TreeViewItem} in the tree, or @code{NULL}
if the tree is empty.
@end deffn



@deffn {GetSelection} gGetSelection::TreeView
@findex gGetSelection::TreeView
@sp 2
@example
@group
tvi = gGetSelection(tv);

object  tv;    /*  TreeView object               */
object  tvi;   /*  selected TreeViewItem or NULL  */
@end group
@end example
Returns the currently selected @code{TreeViewItem}, or @code{NULL} if
no item is selected.
@end deffn



@deffn {LoadBitmap} gLoadBitmap::TreeView
@findex gLoadBitmap::TreeView
@sp 2
@example
@group
id = gLoadBitmap(tv, resID);

object    tv;      /*  TreeView object      */
unsigned  resID;   /*  bitmap resource ID    */
int       id;      /*  image ID or -1       */
@end group
@end example
Loads a bitmap resource into the tree view's image list.  Returns the
image ID for use with @code{gAddTVItemWithImage}, or @code{-1} if no
image list is associated.
@sp 1
See also:  @code{NewImageList}, @code{SetImageList}
@end deffn



@deffn {New} New::TreeView
@findex vNew::TreeView
@sp 2
@example
@group
tv = vNew(TreeView, ctlID, name, dlg);

object    tv;      /*  new TreeView object   */
UINT      ctlID;   /*  control ID            */
char      *name;   /*  name of control       */
object    dlg;     /*  dialog object          */
@end group
@end example
This method creates a new @code{TreeView} control and associates it with
a control defined in the dialog resource identified by @code{ctlID}.
@end deffn



@deffn {NewImageList} gNewImageList::TreeView
@findex gNewImageList::TreeView
@sp 2
@example
@group
r = gNewImageList(tv, cx, cy, flags, numb, grow);

object    tv;      /*  TreeView object      */
int       cx;      /*  image width           */
int       cy;      /*  image height          */
unsigned  flags;   /*  image list flags      */
int       numb;    /*  initial image count   */
int       grow;    /*  grow increment        */
object    r;       /*  TreeView object       */
@end group
@end example
Creates a new @code{ImageList} and associates it with the tree view.
@sp 1
See also:  @code{SetImageList}, @code{LoadBitmap}
@end deffn



@deffn {RemoveAll} gRemoveAll::TreeView
@findex gRemoveAll::TreeView
@sp 2
@example
@group
r = gRemoveAll(tv);

object  tv;   /*  TreeView object  */
object  r;    /*  TreeView object  */
@end group
@end example
Clears the entire tree, removing and disposing all items.
@end deffn



@deffn {Resort} gResort::TreeView
@findex gResort::TreeView
@sp 2
@example
@group
r = gResort(tv);

object  tv;   /*  TreeView object  */
object  r;    /*  TreeView object  */
@end group
@end example
Re-sorts the tree items.  If alphabetical sorting is enabled (via
@code{gAlphabetize}), items are sorted alphabetically.  Otherwise items
are sorted by their @code{order} and @code{suborder} properties.  The
current selection is preserved.
@sp 1
See also:  @code{Alphabetize}, @code{AlphabetizeOff}
@end deffn



@deffn {SetExpFun} gSetExpFun::TreeView
@findex gSetExpFun::TreeView
@sp 2
@example
@group
prev = gSetExpFun(tv, fun);

object  tv;     /*  TreeView object                        */
ifun    fun;    /*  expand/collapse callback function       */
ifun    prev;   /*  previous callback function              */
@end group
@end example
Sets a callback function that is called when a tree item is expanded or
collapsed.  Returns the previous callback function.
@end deffn



@deffn {SetImageList} gSetImageList::TreeView
@findex gSetImageList::TreeView
@sp 2
@example
@group
r = gSetImageList(tv, imageList);

object  tv;          /*  TreeView object     */
object  imageList;   /*  ImageList object    */
object  r;           /*  TreeView object     */
@end group
@end example
Associates an @code{ImageList} object with the tree view for custom item
icons.
@sp 1
See also:  @code{NewImageList}, @code{LoadBitmap}
@end deffn



@deffn {SetSelFun} gSetSelFun::TreeView
@findex gSetSelFun::TreeView
@sp 2
@example
@group
prev = gSetSelFun(tv, fun);

object  tv;     /*  TreeView object                        */
ifun    fun;    /*  selection changed callback function     */
ifun    prev;   /*  previous callback function              */
@end group
@end example
Sets a callback function that is called when the selected tree item
changes.  Returns the previous callback function.  The callback receives
the dialog object, the TreeView control, and the newly selected
@code{TreeViewItem}.
@end deffn



@subsection TreeViewItem Methods



@deffn {AutoDisposeTag} gAutoDisposeTag::TreeViewItem
@findex gAutoDisposeTag::TreeViewItem
@sp 2
@example
@group
r = gAutoDisposeTag(tvi);

object  tvi;   /*  TreeViewItem object  */
object  r;     /*  TreeViewItem object  */
@end group
@end example
Enables automatic disposal of the tag object when the tree item is
disposed or when a new tag replaces the old one.
@sp 1
See also:  @code{SetTag}, @code{GetTag}
@end deffn



@deffn {ChangeImage} gChangeImage::TreeViewItem
@findex gChangeImage::TreeViewItem
@sp 2
@example
@group
r = gChangeImage(tvi, image);

object  tvi;     /*  TreeViewItem object  */
int     image;   /*  new image index      */
object  r;       /*  TreeViewItem object  */
@end group
@end example
Changes the display image of the tree item.  The image index refers to
an image in the @code{ImageList} associated with the parent
@code{TreeView}.
@sp 1
See also:  @code{GetImageID}
@end deffn



@deffn {Child} gChild::TreeViewItem
@findex gChild::TreeViewItem
@sp 2
@example
@group
c = gChild(tvi);

object  tvi;   /*  TreeViewItem object             */
object  c;     /*  first child item or NULL         */
@end group
@end example
Returns the first child @code{TreeViewItem}, or @code{NULL} if the item
has no children.
@end deffn



@deffn {Children} gChildren::TreeViewItem
@findex gChildren::TreeViewItem
@sp 2
@example
@group
c = gChildren(tvi);

object  tvi;   /*  TreeViewItem object              */
object  c;     /*  children collection or NULL       */
@end group
@end example
Returns the collection of child @code{TreeViewItem} objects, or
@code{NULL} if the item has no children.
@end deffn



@deffn {Collapse} gCollapse::TreeViewItem
@findex gCollapse::TreeViewItem
@sp 2
@example
@group
r = gCollapse(tvi);

object  tvi;   /*  TreeViewItem object  */
object  r;     /*  TreeViewItem object  */
@end group
@end example
Collapses the tree node to hide its children.
@sp 1
See also:  @code{Expand}
@end deffn



@deffn {Control} gControl::TreeViewItem
@findex gControl::TreeViewItem
@sp 2
@example
@group
tv = gControl(tvi);

object  tvi;   /*  TreeViewItem object  */
object  tv;    /*  TreeView object      */
@end group
@end example
Returns the @code{TreeView} control object that contains this item.
@end deffn



@deffn {DeleteTVItem} gDeleteTVItem::TreeViewItem
@findex gDeleteTVItem
@sp 2
@example
@group
r = gDeleteTVItem(tvi);

object  tvi;   /*  TreeViewItem object  */
object  r;     /*  NULL                 */
@end group
@end example
Removes this item from its parent and disposes it.  All children are
also disposed.
@end deffn



@deffn {DisposeChildren} gDisposeChildren::TreeViewItem
@findex gDisposeChildren::TreeViewItem
@sp 2
@example
@group
r = gDisposeChildren(tvi);

object  tvi;   /*  TreeViewItem object  */
object  r;
@end group
@end example
Disposes all child items of this node.
@end deffn



@deffn {Expand} gExpand::TreeViewItem
@findex gExpand::TreeViewItem
@sp 2
@example
@group
r = gExpand(tvi);

object  tvi;   /*  TreeViewItem object  */
object  r;     /*  TreeViewItem object  */
@end group
@end example
Expands the tree node to show its children.
@sp 1
See also:  @code{Collapse}
@end deffn



@deffn {GetImageID} gGetImageID::TreeViewItem
@findex gGetImageID::TreeViewItem
@sp 2
@example
@group
id = gGetImageID(tvi);

object  tvi;   /*  TreeViewItem object  */
int     id;    /*  image index          */
@end group
@end example
Returns the current image index for this tree item.
@sp 1
See also:  @code{ChangeImage}
@end deffn



@deffn {GetParent} gGetParent::TreeViewItem
@findex gGetParent::TreeViewItem
@sp 2
@example
@group
p = gGetParent(tvi);

object  tvi;   /*  TreeViewItem object          */
object  p;     /*  parent item or NULL           */
@end group
@end example
Returns the parent @code{TreeViewItem}, or @code{NULL} if the item is
at the root level.
@end deffn



@deffn {GetTag} gGetTag::TreeViewItem
@findex gGetTag::TreeViewItem
@sp 2
@example
@group
tag = gGetTag(tvi);

object  tvi;   /*  TreeViewItem object     */
object  tag;   /*  tag object or NULL      */
@end group
@end example
Returns the user data object attached to the tree item, or @code{NULL}
if none is set.
@sp 1
See also:  @code{SetTag}
@end deffn



@deffn {IsChildOf} gIsChildOf::TreeViewItem
@findex gIsChildOf::TreeViewItem
@sp 2
@example
@group
val = gIsChildOf(tvi, parent);

object  tvi;      /*  TreeViewItem object  */
object  parent;   /*  potential ancestor    */
int     val;      /*  1=yes, 0=no          */
@end group
@end example
Returns 1 if the item is a descendant of @code{parent}, 0 otherwise.
Walks up the parent chain to determine the relationship.
@end deffn



@deffn {NewTVItem} gNewTVItem::TreeViewItem
@findex gNewTVItem
@sp 2
@example
@group
tvi = gNewTVItem(TreeViewItem, tv, parent, image, alpha, txt);

object  tvi;      /*  new TreeViewItem object       */
object  tv;       /*  TreeView control object        */
object  parent;   /*  parent item or NULL             */
int     image;    /*  image index                    */
int     alpha;    /*  alphabetize flag               */
char    *txt;     /*  item display text              */
@end group
@end example
Low-level factory method to create a @code{TreeViewItem}.  Most
applications should use @code{gAddTVItem} or @code{gAddTVItemWithImage}
on the @code{TreeView} instead.
@end deffn



@deffn {Next} gNext::TreeViewItem
@findex gNext::TreeViewItem
@sp 2
@example
@group
nxt = gNext(tvi);

object  tvi;   /*  TreeViewItem object        */
object  nxt;   /*  next item or NULL           */
@end group
@end example
Returns the next visible @code{TreeViewItem} in tree traversal order,
or @code{NULL} if there is no next item.
@end deffn



@deffn {Select} gSelect::TreeViewItem
@findex gSelect::TreeViewItem
@sp 2
@example
@group
r = gSelect(tvi);

object  tvi;   /*  TreeViewItem object  */
object  r;     /*  TreeViewItem object  */
@end group
@end example
Programmatically selects this item in the tree view.
@end deffn



@deffn {SetStringValue} gSetStringValue::TreeViewItem
@findex gSetStringValue::TreeViewItem
@sp 2
@example
@group
r = gSetStringValue(tvi, val);

object  tvi;    /*  TreeViewItem object  */
char    *val;   /*  new display text     */
object  r;      /*  TreeViewItem object  */
@end group
@end example
Updates the display text of the tree item.
@sp 1
See also:  @code{StringValue}
@end deffn



@deffn {SetTag} gSetTag::TreeViewItem
@findex gSetTag::TreeViewItem
@sp 2
@example
@group
prev = gSetTag(tvi, tag);

object  tvi;    /*  TreeViewItem object       */
object  tag;    /*  user data object          */
object  prev;   /*  previous tag or NULL      */
@end group
@end example
Attaches an arbitrary user data object to the tree item.  Returns the
previous tag.  If auto-dispose is enabled and a previous tag exists, it
is automatically disposed.
@sp 1
See also:  @code{GetTag}, @code{AutoDisposeTag}
@end deffn



@deffn {StringValue} gStringValue::TreeViewItem
@findex gStringValue::TreeViewItem
@sp 2
@example
@group
s = gStringValue(tvi);

object  tvi;   /*  TreeViewItem object  */
char    *s;    /*  display text string  */
@end group
@end example
Returns the current display text of the tree item.
@sp 1
See also:  @code{SetStringValue}
@end deffn




@section Tool Bar
The @code{ToolBar} class (a subclass of @code{ChildWindow}) implements a
toolbar control that displays a row of bitmap buttons and optional combo
box controls at the top of a window.  Each button can have an associated
callback function and tooltip.

@subsection ToolBar Methods



@deffn {AddToolBitmap} gAddToolBitmap::ToolBar
@findex gAddToolBitmap
@sp 2
@example
@group
r = gAddToolBitmap(tb, id1, id2, space, fun, tip);

object    tb;      /*  ToolBar object                     */
unsigned  id1;     /*  resource ID of normal bitmap        */
unsigned  id2;     /*  resource ID of pressed bitmap (0=same)  */
int       space;   /*  extra space to the left (pixels)    */
LRESULT   (*fun)();  /*  click callback function           */
char      *tip;    /*  tooltip text or NULL                */
object    r;       /*  ToolBar object                     */
@end group
@end example
Adds a bitmap button to the toolbar.  @code{id1} is the resource ID of
the normal (up) state bitmap.  @code{id2} is the pressed (down) state
bitmap; pass 0 to use the same bitmap for both states.  @code{space}
adds extra horizontal spacing before the button.  @code{fun} is the
callback function invoked when the button is clicked.  @code{tip} is
the tooltip text displayed on hover.

The callback function receives two arguments: the parent window object
and the button resource ID.
@sp 1
See also:  @code{AddToolBitmapFromFile}
@end deffn



@deffn {AddToolBitmapFromFile} gAddToolBitmapFromFile::ToolBar
@findex gAddToolBitmapFromFile
@sp 2
@example
@group
r = gAddToolBitmapFromFile(tb, file, id1, id2, space, fun, tip);

object    tb;      /*  ToolBar object                     */
char      *file;   /*  bitmap file path                    */
unsigned  id1;     /*  resource ID of normal bitmap        */
unsigned  id2;     /*  resource ID of pressed bitmap       */
int       space;   /*  extra space to the left (pixels)    */
LRESULT   (*fun)();  /*  click callback function           */
char      *tip;    /*  tooltip text or NULL                */
object    r;       /*  ToolBar object                     */
@end group
@end example
Adds a bitmap button from an external file rather than a resource.
Otherwise identical to @code{gAddToolBitmap}.
@sp 1
See also:  @code{AddToolBitmap}
@end deffn



@deffn {AddToolComboBox} gAddToolComboBox::ToolBar
@findex gAddToolComboBox
@sp 2
@example
@group
ctl = gAddToolComboBox(tb, height, width, space, fun, tip, name);

object  tb;      /*  ToolBar object                    */
int     height;  /*  drop-down height                   */
int     width;   /*  control width                      */
int     space;   /*  extra space to the left (pixels)   */
int     (*fun)();  /*  change callback function          */
char    *tip;    /*  tooltip text or NULL                */
char    *name;   /*  control name                       */
object  ctl;     /*  ComboBox control object             */
@end group
@end example
Adds a @code{ComboBox} control to the toolbar.  Returns the combo box
control object, which can be configured using the @code{ComboBox}
methods.
@end deffn



@deffn {ChangeTip} gChangeTip::ToolBar
@findex gChangeTip
@sp 2
@example
@group
r = gChangeTip(tb, id, tip);

object    tb;    /*  ToolBar object            */
unsigned  id;    /*  button resource ID         */
char      *tip;  /*  new tooltip text           */
object    r;     /*  ToolBar object             */
@end group
@end example
Changes the tooltip text for a toolbar button.
@sp 1
See also:  @code{GetTip}
@end deffn



@deffn {EnableToolBitmap} gEnableToolBitmap::ToolBar
@findex gEnableToolBitmap
@sp 2
@example
@group
changed = gEnableToolBitmap(tb, id, val);

object    tb;       /*  ToolBar object              */
unsigned  id;       /*  button resource ID           */
int       val;      /*  1=enable, 0=disable          */
int       changed;  /*  1 if state changed, 0 if not */
@end group
@end example
Enables or disables a toolbar button.  Disabled buttons are drawn with
a grayed-out appearance and do not respond to clicks.
@end deffn



@deffn {GetControlStr} gGetControlStr::ToolBar
@findex gGetControlStr::ToolBar
@sp 2
@example
@group
ctl = gGetControlStr(tb, name);

object  tb;     /*  ToolBar object                     */
char    *name;  /*  control name                        */
object  ctl;    /*  ComboBox control object or NULL     */
@end group
@end example
Retrieves a combo box control embedded in the toolbar by its name.
Returns @code{NULL} if not found.
@end deffn



@deffn {GetTip} gGetTip::ToolBar
@findex gGetTip
@sp 2
@example
@group
s = gGetTip(tb, id);

object    tb;   /*  ToolBar object       */
unsigned  id;   /*  button resource ID    */
char      *s;   /*  tooltip text          */
@end group
@end example
Returns the tooltip text for a toolbar button.
@sp 1
See also:  @code{ChangeTip}
@end deffn



@deffn {New} New::ToolBar
@findex vNew::ToolBar
@sp 2
@example
@group
tb = vNew(ToolBar, win);

object  tb;    /*  new ToolBar object    */
object  win;   /*  parent window object  */
@end group
@end example
Creates a new @code{ToolBar} and attaches it to the top of the specified
parent window.
@end deffn



@deffn {PixelHeight} gPixelHeight::ToolBar
@findex gPixelHeight::ToolBar
@sp 2
@example
@group
h = gPixelHeight(tb);

object  tb;   /*  ToolBar object         */
int     h;    /*  toolbar pixel height    */
@end group
@end example
Returns the pixel height of the toolbar.
@end deffn



@deffn {SetAccessModeFunction} gSetAccessModeFunction::ToolBar
@findex gSetAccessModeFunction::ToolBar
@sp 2
@example
@group
prev = gSetAccessModeFunction(ToolBar, fun);

ifun  fun;    /*  access mode function  */
ifun  prev;   /*  previous function     */
@end group
@end example
This is a class method.  Sets a function that determines whether each
toolbar button should be enabled or disabled based on application-level
access control.  Returns the previous function.

The function is called for each button during @code{gUpdateAccessMode}
and receives the button resource ID as an argument.
@sp 1
See also:  @code{UpdateAccessMode}
@end deffn



@deffn {SetToolBarMouseFunction} gSetToolBarMouseFunction::ToolBar
@findex gSetToolBarMouseFunction
@sp 2
@example
@group
prev = gSetToolBarMouseFunction(tb, id, button, fun);

object    tb;       /*  ToolBar object                       */
unsigned  id;       /*  button resource ID                    */
unsigned  button;   /*  mouse button/modifier combination     */
ifun      fun;      /*  callback function                     */
ofun      prev;     /*  previous callback function            */
@end group
@end example
Sets a mouse event callback for a specific toolbar button and
mouse button/modifier combination.  The @code{button} parameter
encodes both the mouse button (left/right) and modifier keys
(Shift, Control).
@end deffn



@deffn {SetToolFunction} gSetToolFunction::ToolBar
@findex gSetToolFunction
@sp 2
@example
@group
prev = gSetToolFunction(tb, id, fun);

object    tb;     /*  ToolBar object                */
unsigned  id;     /*  button resource ID             */
LRESULT   (*fun)();  /*  new click callback          */
ifun      prev;   /*  previous callback function     */
@end group
@end example
Changes the click callback function for a toolbar button.  Returns the
previous function.
@end deffn



@deffn {ShowToolBitmap} gShowToolBitmap::ToolBar
@findex gShowToolBitmap
@sp 2
@example
@group
prev = gShowToolBitmap(tb, id, val);

object    tb;     /*  ToolBar object           */
unsigned  id;     /*  button resource ID        */
int       val;    /*  1=show, 0=hide            */
int       prev;   /*  previous hidden state     */
@end group
@end example
Shows or hides a toolbar button.
@end deffn



@deffn {UpdateAccessMode} gUpdateAccessMode::ToolBar
@findex gUpdateAccessMode
@sp 2
@example
@group
r = gUpdateAccessMode(tb);

object  tb;   /*  ToolBar object  */
object  r;    /*  ToolBar object  */
@end group
@end example
Updates the enabled/disabled state of all toolbar buttons by calling
the access mode function set via @code{gSetAccessModeFunction}.
@sp 1
See also:  @code{SetAccessModeFunction}
@end deffn



@deffn {UpdateState} gUpdateState::ToolBar
@findex gUpdateState::ToolBar
@sp 2
@example
@group
r = gUpdateState(tb);

object  tb;   /*  ToolBar object  */
object  r;    /*  ToolBar object  */
@end group
@end example
Updates the toolbar's position and size to match the parent window's
current width.  Typically called after the parent window is resized.
@end deffn



@deffn {UseAlternateColor} gUseAlternateColor::ToolBar
@findex gUseAlternateColor
@sp 2
@example
@group
gUseAlternateColor(tb);

object  tb;   /*  ToolBar object  */
@end group
@end example
Enables alternate color mode for toolbar button rendering.
@end deffn




@section Status Bar
The @code{StatusBar} class (a subclass of @code{ChildWindow}) implements
a status bar control that displays one or more text sections at the
bottom of a window.

@subsection StatusBar Methods



@deffn {AddSection} gAddSection::StatusBar
@findex gAddSection
@sp 2
@example
@group
r = gAddSection(sb, id, width, mode);

object  sb;      /*  StatusBar object           */
int     id;      /*  section identifier          */
int     width;   /*  section width               */
int     mode;    /*  alignment mode              */
object  r;       /*  StatusBar object or NULL    */
@end group
@end example
Adds a new section to the status bar.  @code{id} is used to identify the
section when setting text.  @code{width} specifies the section width.
Returns @code{NULL} if the maximum number of sections (20) has been
reached.
@sp 1
See also:  @code{SetSectionText}
@end deffn



@deffn {New} New::StatusBar
@findex vNew::StatusBar
@sp 2
@example
@group
sb = vNew(StatusBar, win);

object  sb;    /*  new StatusBar object   */
object  win;   /*  parent window object   */
@end group
@end example
Creates a new @code{StatusBar} and attaches it to the bottom of the
specified parent window.
@end deffn



@deffn {PixelHeight} gPixelHeight::StatusBar
@findex gPixelHeight::StatusBar
@sp 2
@example
@group
h = gPixelHeight(sb);

object  sb;   /*  StatusBar object        */
int     h;    /*  status bar pixel height  */
@end group
@end example
Returns the pixel height of the status bar.
@end deffn



@deffn {SetSectionText} gSetSectionText::StatusBar
@findex gSetSectionText
@sp 2
@example
@group
r = gSetSectionText(sb, id, text);

object  sb;     /*  StatusBar object  */
int     id;     /*  section ID        */
char    *text;  /*  display text      */
object  r;      /*  StatusBar object  */
@end group
@end example
Sets or updates the text displayed in a status bar section.
@sp 1
See also:  @code{AddSection}
@end deffn



@deffn {UpdateState} gUpdateState::StatusBar
@findex gUpdateState::StatusBar
@sp 2
@example
@group
r = gUpdateState(sb);

object  sb;   /*  StatusBar object  */
object  r;    /*  StatusBar object  */
@end group
@end example
Updates the status bar's position and size to align with the parent
window's bottom edge.  Typically called after the parent window is
resized.
@end deffn




@section Image Control
The @code{ImageControl} class (a subclass of @code{Control}) implements
a control for displaying bitmap images in dialogs and windows.  Bitmaps
can be loaded from resource IDs or from external @code{.bmp} files.

@subsection ImageControl Methods



@deffn {GetExitType} gGetExitType::ImageControl
@findex gGetExitType::ImageControl
@sp 2
@example
@group
type = gGetExitType(ic);

object  ic;     /*  ImageControl object  */
int     type;   /*  exit type            */
@end group
@end example
Returns the exit type previously set by @code{gSetExitType}.
@sp 1
See also:  @code{SetExitType}
@end deffn



@deffn {New} New::ImageControl
@findex vNew::ImageControl
@sp 2
@example
@group
ic = vNew(ImageControl, ctlID, name, dlg);

object  ic;       /*  new ImageControl object  */
UINT    ctlID;    /*  control ID               */
char    *name;    /*  control name             */
object  dlg;      /*  dialog object            */
@end group
@end example
Creates a new @code{ImageControl} associated with a control defined in
the dialog resource.
@end deffn



@deffn {SetBitmapID} gSetBitmapID::ImageControl
@findex gSetBitmapID
@sp 2
@example
@group
r = gSetBitmapID(ic, id);

object  ic;   /*  ImageControl object       */
int     id;   /*  bitmap resource ID         */
object  r;    /*  ImageControl object        */
@end group
@end example
Sets the displayed bitmap from a resource ID.  The bitmap is loaded and
displayed when the control is shown.
@sp 1
See also:  @code{SetStringValue}
@end deffn



@deffn {SetButtonBorder} gSetButtonBorder::ImageControl
@findex gSetButtonBorder
@sp 2
@example
@group
prev = gSetButtonBorder(ic, val);

object  ic;     /*  ImageControl object              */
int     val;    /*  -1=down, 0=none, 1=up border     */
int     prev;   /*  previous border value             */
@end group
@end example
Sets the 3D border style around the image control.  @code{1} gives a
raised border, @code{-1} gives a sunken border, and @code{0} removes
the border.  Returns the previous border setting.
@end deffn



@deffn {SetDCFunction} gSetDCFunction::ImageControl
@findex gSetDCFunction::ImageControl
@sp 2
@example
@group
prev = gSetDCFunction(ic, fun);

object  ic;     /*  ImageControl object               */
int     (*fun)();  /*  double-click callback function  */
ofun    prev;   /*  previous callback function         */
@end group
@end example
Sets a callback function that is invoked when the user double-clicks
the image control.  Returns the previous callback function.
@end deffn



@deffn {SetExitType} gSetExitType::ImageControl
@findex gSetExitType::ImageControl
@sp 2
@example
@group
r = gSetExitType(ic, type);

object  ic;     /*  ImageControl object        */
int     type;   /*  IDOK, IDCANCEL, or 0       */
object  r;      /*  ImageControl object        */
@end group
@end example
Configures the image control to act like a button that closes the dialog
when clicked.  Pass @code{IDOK} or @code{IDCANCEL} to enable this
behavior, or 0 to disable it.
@sp 1
See also:  @code{GetExitType}
@end deffn



@deffn {SetStringValue} gSetStringValue::ImageControl
@findex gSetStringValue::ImageControl
@sp 2
@example
@group
r = gSetStringValue(ic, val);

object  ic;    /*  ImageControl object   */
char    *val;  /*  bitmap file path      */
object  r;     /*  ImageControl object   */
@end group
@end example
Sets the displayed bitmap from an external @code{.bmp} file path.
@sp 1
See also:  @code{SetBitmapID}, @code{StringValue}
@end deffn



@deffn {StringValue} gStringValue::ImageControl
@findex gStringValue::ImageControl
@sp 2
@example
@group
s = gStringValue(ic);

object  ic;   /*  ImageControl object   */
char    *s;   /*  bitmap file path      */
@end group
@end example
Returns the bitmap file path currently associated with the control.
@sp 1
See also:  @code{SetStringValue}
@end deffn




@section Time Control
The @code{TimeControl} class (a subclass of @code{SpinControl}) implements
a control for entering and validating time values.  It supports
configurable time ranges, formatting, and 12-hour wrapping.

@subsection TimeControl Methods



@deffn {BlankIfZero} gBlankIfZero::TimeControl
@findex gBlankIfZero
@sp 2
@example
@group
prev = gBlankIfZero(tc, val);

object  tc;     /*  TimeControl object     */
int     val;    /*  1=blank if zero, 0=no  */
int     prev;   /*  previous setting       */
@end group
@end example
When enabled, the control displays as blank when the time value is zero
rather than showing @code{0:00}.  Returns the previous setting.
@end deffn



@deffn {FormatTime} gFormatTime::TimeControl
@findex gFormatTime
@sp 2
@example
@group
r = gFormatTime(tc, fmt);

object  tc;    /*  TimeControl object     */
char    *fmt;  /*  format string          */
object  r;     /*  TimeControl object     */
@end group
@end example
Sets the format string used to display the time value.
@end deffn



@deffn {HideData} gHideData::TimeControl
@findex gHideData::TimeControl
@sp 2
@example
@group
prev = gHideData(tc, flg);

object  tc;     /*  TimeControl object        */
int     flg;    /*  1=hide, 0=show            */
int     prev;   /*  previous setting          */
@end group
@end example
When enabled, the control displays asterisks instead of the actual time
value.  Returns the previous setting.
@end deffn



@deffn {IsBlank} gIsBlank::TimeControl
@findex gIsBlank::TimeControl
@sp 2
@example
@group
val = gIsBlank(tc);

object  tc;    /*  TimeControl object     */
int     val;   /*  1=blank, 0=has value   */
@end group
@end example
Returns 1 if the control is currently blank, 0 otherwise.
@end deffn



@deffn {New} New::TimeControl
@findex vNew::TimeControl
@sp 2
@example
@group
tc = vNew(TimeControl, ctlID, name, dlg);

object  tc;       /*  new TimeControl object  */
UINT    ctlID;    /*  control ID              */
char    *name;    /*  control name            */
object  dlg;      /*  dialog object           */
@end group
@end example
Creates a new @code{TimeControl} associated with a control defined in
the dialog resource.
@end deffn



@deffn {SetBeginHour} gSetBeginHour::TimeControl
@findex gSetBeginHour::TimeControl
@sp 2
@example
@group
prev = gSetBeginHour(tc, hour);

object  tc;     /*  TimeControl object        */
int     hour;   /*  beginning hour for wrap   */
int     prev;   /*  previous begin hour       */
@end group
@end example
Sets the beginning hour for 12-hour wrap-around behavior.  For example,
setting this to 8 means hours are interpreted in the range 8:00--7:59.
Returns the previous value.

This method may also be called as a class method on @code{TimeControl}
to set the default for all instances.
@end deffn



@deffn {SetDefaultLong} gSetDefaultLong::TimeControl
@findex gSetDefaultLong::TimeControl
@sp 2
@example
@group
r = gSetDefaultLong(tc, val);

object  tc;    /*  TimeControl object  */
long    val;   /*  default time value  */
object  r;     /*  TimeControl object  */
@end group
@end example
Sets the default time value used when @code{gUseDefault} is called.
@end deffn



@deffn {SetTimeValue} gSetTimeValue::TimeControl
@findex gSetTimeValue
@sp 2
@example
@group
r = gSetTimeValue(tc, val);

object  tc;    /*  TimeControl object  */
long    val;   /*  time as long value  */
object  r;     /*  TimeControl object  */
@end group
@end example
Sets the time value from a long integer and updates the display.
@code{gSetLongValue} is a synonym for this method.
@sp 1
See also:  @code{TimeValue}
@end deffn



@deffn {TimeRange} gTimeRange::TimeControl
@findex gTimeRange
@sp 2
@example
@group
r = gTimeRange(tc, minimum, maximum, allowNone);

object  tc;          /*  TimeControl object          */
long    minimum;     /*  minimum time value           */
long    maximum;     /*  maximum time value           */
int     allowNone;   /*  1=allow empty, 0=required    */
object  r;           /*  TimeControl object           */
@end group
@end example
Sets the valid time range for the control.  If @code{allowNone} is 1,
the user may leave the field empty.
@end deffn



@deffn {TimeValue} gTimeValue::TimeControl
@findex gTimeValue
@sp 2
@example
@group
val = gTimeValue(tc);

object  tc;    /*  TimeControl object  */
long    val;   /*  time as long value  */
@end group
@end example
Returns the current time value as a long integer.
@code{gLongValue} is a synonym for this method.
@sp 1
See also:  @code{SetTimeValue}
@end deffn




@section Shortcut Menu
The @code{ShortcutMenu} class (a subclass of @code{Menu}) implements
context menus (also known as popup or right-click menus) that can be
displayed at arbitrary screen positions.

@subsection ShortcutMenu Methods



@deffn {AddMenuOption} gAddMenuOption::ShortcutMenu
@findex gAddMenuOption::ShortcutMenu
@sp 2
@example
@group
id = gAddMenuOption(sm, title, fun);

object  sm;      /*  ShortcutMenu object        */
char    *title;  /*  menu item text              */
void    (*fun)();  /*  callback function         */
int     id;      /*  assigned menu item ID       */
@end group
@end example
Dynamically adds a menu item with the given text and callback function.
Returns the automatically assigned ID for the menu item.
@sp 1
See also:  @code{AddSubMenu}, @code{AddSeparator}
@end deffn



@deffn {AddSeparator} gAddSeparator::ShortcutMenu
@findex gAddSeparator::ShortcutMenu
@sp 2
@example
@group
r = gAddSeparator(sm);

object  sm;   /*  ShortcutMenu object  */
object  r;    /*  ShortcutMenu object  */
@end group
@end example
Adds a separator line to the menu.
@sp 1
See also:  @code{AddMenuOption}
@end deffn



@deffn {AddSubMenu} gAddSubMenu::ShortcutMenu
@findex gAddSubMenu
@sp 2
@example
@group
sub = gAddSubMenu(sm, title);

object  sm;      /*  ShortcutMenu object          */
char    *title;  /*  submenu title                 */
object  sub;     /*  new submenu object            */
@end group
@end example
Creates and adds a new submenu with the given title.  Returns the new
submenu object, which can have items added to it via
@code{gAddMenuOption}.
@sp 1
See also:  @code{AddMenuOption}, @code{AddSeparator}
@end deffn



@deffn {LoadShortcutMenu} gLoadShortcutMenu::ShortcutMenu
@findex gLoadShortcutMenu
@sp 2
@example
@group
sm = gLoadShortcutMenu(ShortcutMenu, id, submenu);

object    sm;        /*  new ShortcutMenu object   */
unsigned  id;        /*  menu resource ID           */
int       submenu;   /*  submenu index              */
@end group
@end example
This is a class method.  Loads a shortcut menu from a menu resource by
numeric ID.  @code{submenu} specifies which submenu to use from the
resource.
@sp 1
See also:  @code{LoadShortcutMenuStr}
@end deffn



@deffn {LoadShortcutMenuFromFile} gLoadShortcutMenuFromFile::ShortcutMenu
@findex gLoadShortcutMenuFromFile
@sp 2
@example
@group
sm = gLoadShortcutMenuFromFile(ShortcutMenu, file, id, submenu);

object    sm;        /*  new ShortcutMenu object   */
char      *file;     /*  resource file path         */
unsigned  id;        /*  menu resource ID           */
int       submenu;   /*  submenu index              */
@end group
@end example
This is a class method.  Loads a shortcut menu from an external resource
file.
@end deffn



@deffn {LoadShortcutMenuStr} gLoadShortcutMenuStr::ShortcutMenu
@findex gLoadShortcutMenuStr
@sp 2
@example
@group
sm = gLoadShortcutMenuStr(ShortcutMenu, name, submenu);

object  sm;        /*  new ShortcutMenu object   */
char    *name;     /*  menu resource name         */
int     submenu;   /*  submenu index              */
@end group
@end example
This is a class method.  Loads a shortcut menu from a menu resource by
string name.
@sp 1
See also:  @code{LoadShortcutMenu}
@end deffn



@deffn {New} New::ShortcutMenu
@findex vNew::ShortcutMenu
@sp 2
@example
@group
sm = vNew(ShortcutMenu);

object  sm;   /*  new ShortcutMenu object  */
@end group
@end example
Creates a new, empty shortcut menu.
@sp 1
See also:  @code{LoadShortcutMenu}, @code{LoadShortcutMenuStr}
@end deffn



@deffn {Perform} gPerform::ShortcutMenu
@findex gPerform::ShortcutMenu
@sp 2
@example
@group
r = gPerform(sm);

object  sm;   /*  ShortcutMenu object  */
int     r;    /*  result               */
@end group
@end example
Displays the shortcut menu at the current cursor position and executes
the callback associated with the item selected by the user.
@end deffn




@section Static Control
The @code{StaticControl} class (a subclass of @code{Control}) implements
a basic static text display control.  The @code{StaticTextControl} class
(a subclass of @code{TextControl}) provides additional capabilities
including printing, XML serialization, and multi-language support.

@subsection StaticControl Methods



@deffn {New} New::StaticControl
@findex vNew::StaticControl
@sp 2
@example
@group
sc = vNew(StaticControl, ctlID, name, dlg);

object  sc;       /*  new StaticControl object  */
UINT    ctlID;    /*  control ID                */
char    *name;    /*  control name              */
object  dlg;      /*  dialog object             */
@end group
@end example
Creates a new @code{StaticControl} associated with a control defined in
the dialog resource.
@end deffn



@deffn {SetStringValue} gSetStringValue::StaticControl
@findex gSetStringValue::StaticControl
@sp 2
@example
@group
r = gSetStringValue(sc, val);

object  sc;    /*  StaticControl object  */
char    *val;  /*  display text          */
object  r;     /*  StaticControl object  */
@end group
@end example
Sets the text displayed by the static control.  @code{gSetTitle} is a
synonym for this method.
@sp 1
See also:  @code{StringValue}
@end deffn



@deffn {StringValue} gStringValue::StaticControl
@findex gStringValue::StaticControl
@sp 2
@example
@group
s = gStringValue(sc);

object  sc;   /*  StaticControl object  */
char    *s;   /*  display text string   */
@end group
@end example
Returns the current display text.
@sp 1
See also:  @code{SetStringValue}
@end deffn



@subsection StaticTextControl Methods

The @code{StaticTextControl} class inherits from @code{TextControl} and
adds support for URL bindings, multi-language text, and right-alignment.



@deffn {ResetText} gResetText::StaticTextControl
@findex gResetText
@sp 2
@example
@group
r = gResetText(stc);

object  stc;   /*  StaticTextControl object  */
object  r;     /*  StaticTextControl object  */
@end group
@end example
Resets the window text to the current language's text value.
@sp 1
See also:  @code{SetCurrentLanguage}
@end deffn



@deffn {SetCurrentLanguage} gSetCurrentLanguage::StaticTextControl
@findex gSetCurrentLanguage::StaticTextControl
@sp 2
@example
@group
r = gSetCurrentLanguage(stc, lang);

object  stc;    /*  StaticTextControl object  */
int     lang;   /*  language ID               */
object  r;      /*  StaticTextControl object  */
@end group
@end example
Sets the display language and updates the text accordingly.
@sp 1
See also:  @code{ResetText}
@end deffn



@deffn {SetURL} gSetURL::StaticTextControl
@findex gSetURL
@sp 2
@example
@group
r = gSetURL(stc, url);

object  stc;    /*  StaticTextControl object  */
char    *url;   /*  URL binding string        */
object  r;      /*  StaticTextControl object  */
@end group
@end example
Associates a URL with the static text control.
@sp 1
See also:  @code{URL}
@end deffn



@deffn {URL} gURL::StaticTextControl
@findex gURL
@sp 2
@example
@group
s = gURL(stc);

object  stc;   /*  StaticTextControl object  */
char    *s;    /*  URL string                */
@end group
@end example
Returns the URL associated with this control, or an empty string if
none is set.
@sp 1
See also:  @code{SetURL}
@end deffn




@section Bitmap
The @code{Bitmap} class provides support for displaying bitmap images on
windows.  Bitmaps can be positioned, centered, tiled, or stretched to
fill the window.

@subsection Bitmap Methods



@deffn {DrawBitmap} gDrawBitmap::Bitmap
@findex gDrawBitmap
@sp 2
@example
@group
r = gDrawBitmap(bm);

object  bm;   /*  Bitmap object  */
object  r;    /*  Bitmap object  */
@end group
@end example
Draws the bitmap on its parent window using the positioning mode
specified at creation time (positioned, centered, or tiled).
@end deffn



@deffn {NewBitmap} gNewBitmap::Bitmap
@findex gNewBitmap
@sp 2
@example
@group
bm = gNewBitmap(Bitmap, bmpID, mode, y, x, wind);

object    bm;      /*  new Bitmap object        */
UINT      bmpID;   /*  bitmap resource ID        */
int       mode;    /*  display mode              */
int       y;       /*  vertical position          */
int       x;       /*  horizontal position        */
object    wind;    /*  parent window object       */
@end group
@end example
This is a class method.  Creates a new @code{Bitmap} object from a
resource.  @code{mode} controls how the bitmap is displayed: 1 for
positioned at (@code{x},@code{y}), 2 for centered, 3 for tiled.
Position values are scaled to font pixels.
@sp 1
See also:  @code{NewBitmap2}, @code{DrawBitmap}
@end deffn



@deffn {NewBitmap2} gNewBitmap2::Bitmap
@findex gNewBitmap2
@sp 2
@example
@group
bm = gNewBitmap2(Bitmap, bmpID, mode, y, x, wind);

object    bm;      /*  new Bitmap object        */
UINT      bmpID;   /*  bitmap resource ID        */
int       mode;    /*  display mode              */
int       y;       /*  vertical position          */
int       x;       /*  horizontal position        */
object    wind;    /*  parent window object       */
@end group
@end example
This is a class method.  Identical to @code{gNewBitmap} except that
position values are not scaled to font pixels --- they are used as
raw pixel coordinates.
@sp 1
See also:  @code{NewBitmap}
@end deffn




@section Spin Control
The @code{SpinControl} class (a subclass of @code{Control}) provides
spin button (up/down arrow) functionality that can be attached to other
controls such as @code{TimeControl} or @code{NumericControl}.

@subsection SpinControl Methods



@deffn {SetSpinFunction} gSetSpinFunction::SpinControl
@findex gSetSpinFunction
@sp 2
@example
@group
r = gSetSpinFunction(ctl, spinID, func);

object    ctl;      /*  SpinControl object               */
unsigned  spinID;   /*  resource ID of spin button        */
int       (*func)(object ctl, object dlg, int val);
                    /*  callback function                 */
object    r;        /*  SpinControl object                */
@end group
@end example
Associates a spin button control with this control and sets the callback
function.  The callback receives the control object, the dialog object,
and a value of @code{+1} (up arrow) or @code{-1} (down arrow).
@end deffn



@section Task List
The @code{TaskList} class provides a workflow management system that
presents the user with a list of tasks (steps) displayed in a modeless
dialog.  Each task in the list is represented by a @code{Task} object
that manages its own dialog or processing step.  The user selects tasks
from the list, and the system handles transitions, validation, and
cleanup between tasks.

@subsection TaskList Methods



@deffn {AddCompletionFunctionAfter} gAddCompletionFunctionAfter::TaskList
@findex gAddCompletionFunctionAfter::TaskList
@sp 2
@example
@group
r = gAddCompletionFunctionAfter(tl, fun);

object  tl;    /*  TaskList object          */
int     (*fun)();  /*  completion function  */
object  r;     /*  TaskList object          */
@end group
@end example
Adds a completion function that is called after other completion
functions when the task list closes.
@sp 1
See also:  @code{AddCompletionFunctionBefore}, @code{CompletionFunction}
@end deffn



@deffn {AddCompletionFunctionBefore} gAddCompletionFunctionBefore::TaskList
@findex gAddCompletionFunctionBefore::TaskList
@sp 2
@example
@group
r = gAddCompletionFunctionBefore(tl, fun);

object  tl;    /*  TaskList object          */
int     (*fun)();  /*  completion function  */
object  r;     /*  TaskList object          */
@end group
@end example
Adds a completion function that is called before other completion
functions when the task list closes.
@sp 1
See also:  @code{AddCompletionFunctionAfter}, @code{CompletionFunction}
@end deffn



@deffn {AddTask} gAddTask::TaskList
@findex gAddTask
@sp 2
@example
@group
r = gAddTask(tl, proc, option);

object  tl;      /*  TaskList object               */
object  proc;    /*  task procedure                 */
char    *option; /*  display text for list box      */
object  r;       /*  link object                    */
@end group
@end example
Adds a task to the task list with the given display text.  The text
appears in the list box presented to the user.
@end deffn



@deffn {CheckValue} gCheckValue::TaskList
@findex gCheckValue::TaskList
@sp 2
@example
@group
err = gCheckValue(tl);

object  tl;    /*  TaskList object          */
int     err;   /*  0=valid, non-zero=error  */
@end group
@end example
Validates the current task's dialog values.  Returns 0 if valid,
non-zero if validation fails.
@end deffn



@deffn {CompletionFunction} gCompletionFunction::TaskList
@findex gCompletionFunction::TaskList
@sp 2
@example
@group
r = gCompletionFunction(tl, fun);

object  tl;    /*  TaskList object          */
int     (*fun)();  /*  completion function  */
object  r;     /*  TaskList object          */
@end group
@end example
Sets a completion function to be called when the task list closes.
Replaces any previously set completion function.
@sp 1
See also:  @code{AddCompletionFunctionAfter}, @code{AddCompletionFunctionBefore}
@end deffn



@deffn {CurrentTask} gCurrentTask::TaskList
@findex gCurrentTask
@sp 2
@example
@group
task = gCurrentTask(tl);

object  tl;     /*  TaskList object            */
object  task;   /*  current Task or NULL       */
@end group
@end example
Returns the currently active @code{Task} object, or @code{NULL} if no
task is active.
@end deffn



@deffn {Dialog} gDialog::TaskList
@findex gDialog::TaskList
@sp 2
@example
@group
dlg = gDialog(tl);

object  tl;    /*  TaskList object         */
object  dlg;   /*  ModelessDialog object   */
@end group
@end example
Returns the @code{ModelessDialog} object that contains the task list UI.
@end deffn



@deffn {Disable} gDisable::TaskList
@findex gDisable::TaskList
@sp 2
@example
@group
r = gDisable(tl);

object  tl;   /*  TaskList object  */
object  r;    /*  TaskList object  */
@end group
@end example
Disables the task list dialog.
@sp 1
See also:  @code{Enable}
@end deffn



@deffn {Enable} gEnable::TaskList
@findex gEnable::TaskList
@sp 2
@example
@group
r = gEnable(tl);

object  tl;   /*  TaskList object  */
object  r;    /*  TaskList object  */
@end group
@end example
Enables the task list dialog and restores focus to the current task
if one is active.
@sp 1
See also:  @code{Disable}
@end deffn



@deffn {GetPosition} gGetPosition::TaskList
@findex gGetPosition::TaskList
@sp 2
@example
@group
r = gGetPosition(tl, y, x);

object  tl;   /*  TaskList object               */
int     *y;   /*  pointer to vertical position   */
int     *x;   /*  pointer to horizontal position  */
object  r;    /*  TaskList object                */
@end group
@end example
Retrieves the current position of the task list dialog.
@end deffn



@deffn {GetTag} gGetTag::TaskList
@findex gGetTag::TaskList
@sp 2
@example
@group
tag = gGetTag(tl);

object  tl;    /*  TaskList object      */
object  tag;   /*  tag object or NULL   */
@end group
@end example
Returns the user data object attached to the task list, or @code{NULL}
if none is set.
@sp 1
See also:  @code{SetTag}, @code{AutoDisposeTag}
@end deffn



@deffn {NewTaskList} gNewTaskList::TaskList
@findex gNewTaskList
@sp 2
@example
@group
tl = gNewTaskList(TaskList, dlg_id, lb_id, pwind, fun);

object    tl;      /*  new TaskList object             */
unsigned  dlg_id;  /*  dialog resource ID               */
unsigned  lb_id;   /*  list box control ID              */
object    pwind;   /*  parent window                    */
int       (*fun)();  /*  initialization function or NULL */
@end group
@end example
This is a class method.  Creates a new @code{TaskList} with a modeless
dialog (@code{dlg_id}) containing a list box (@code{lb_id}).
@code{pwind} is the parent window.  @code{fun} is an optional
initialization function called when the dialog is set up.
@end deffn



@deffn {NextTask} gNextTask::TaskList
@findex gNextTask::TaskList
@sp 2
@example
@group
r = gNextTask(tl, result);

object  tl;       /*  TaskList object         */
int     result;   /*  task completion result  */
object  r;        /*  TaskList object         */
@end group
@end example
Advances to the next task.  The @code{result} parameter indicates how
the current task completed (e.g.@: OK or Cancel), which determines
whether the current task's save or cleanup phase is executed.
@end deffn



@deffn {Perform} gPerform::TaskList
@findex gPerform::TaskList
@sp 2
@example
@group
r = gPerform(tl);

object  tl;   /*  TaskList object  */
int     r;    /*  0                */
@end group
@end example
Displays the task list selection dialog non-modally.
@sp 1
See also:  @code{PerformModal}
@end deffn



@deffn {PerformModal} gPerformModal::TaskList
@findex gPerformModal::TaskList
@sp 2
@example
@group
r = gPerformModal(tl);

object  tl;   /*  TaskList object  */
int     r;    /*  0                */
@end group
@end example
Displays the task list selection dialog modally, disabling the parent
window until the task list is closed.
@sp 1
See also:  @code{Perform}
@end deffn



@deffn {SetShortValue} gSetShortValue::TaskList
@findex gSetShortValue::TaskList
@sp 2
@example
@group
r = gSetShortValue(tl, val);

object  tl;    /*  TaskList object          */
int     val;   /*  task index to select     */
object  r;     /*  TaskList object          */
@end group
@end example
Selects a task by index in the list box.
@end deffn



@deffn {SetTag} gSetTag::TaskList
@findex gSetTag::TaskList
@sp 2
@example
@group
prev = gSetTag(tl, tag);

object  tl;     /*  TaskList object          */
object  tag;    /*  user data object         */
object  prev;   /*  previous tag or NULL     */
@end group
@end example
Attaches a user data object to the task list.  Returns the previous tag.
If auto-dispose is enabled and a previous tag exists, it is automatically
disposed.
@sp 1
See also:  @code{GetTag}, @code{AutoDisposeTag}
@end deffn



@deffn {SetZOrder} gSetZOrder::TaskList
@findex gSetZOrder::TaskList
@sp 2
@example
@group
r = gSetZOrder(tl, hwnd);

object  tl;     /*  TaskList object           */
HWND    hwnd;   /*  window to position after  */
object  r;      /*  TaskList object           */
@end group
@end example
Sets the z-order (window stacking position) for the task list dialog.
@end deffn



@subsection Task Methods



@deffn {AutoDisposeTag} gAutoDisposeTag::Task
@findex gAutoDisposeTag::Task
@sp 2
@example
@group
r = gAutoDisposeTag(task);

object  task;   /*  Task object  */
object  r;      /*  Task object  */
@end group
@end example
Enables automatic disposal of the tag object when the task is disposed
or when a new tag replaces the old one.
@sp 1
See also:  @code{SetTag}, @code{GetTag}
@end deffn



@deffn {CheckValue} gCheckValue::Task
@findex gCheckValue::Task
@sp 2
@example
@group
err = gCheckValue(task);

object  task;   /*  Task object              */
int     err;    /*  0=valid, non-zero=error  */
@end group
@end example
Validates the task's dialog values.  Returns 0 if valid, non-zero if
validation fails.
@end deffn



@deffn {GetTag} gGetTag::Task
@findex gGetTag::Task
@sp 2
@example
@group
tag = gGetTag(task);

object  task;   /*  Task object          */
object  tag;    /*  tag object or NULL   */
@end group
@end example
Returns the user data object attached to the task, or @code{NULL} if
none is set.
@sp 1
See also:  @code{SetTag}, @code{AutoDisposeTag}
@end deffn



@deffn {NewTask} gNewTask::Task
@findex gNewTask
@sp 2
@example
@group
task = gNewTask(Task, tl);

object  task;   /*  new Task object       */
object  tl;     /*  parent TaskList       */
@end group
@end example
This is a class method.  Creates a new @code{Task} associated with the
given @code{TaskList}.
@end deffn



@deffn {NextTask} gNextTask::Task
@findex gNextTask::Task
@sp 2
@example
@group
r = gNextTask(task, result);

object  task;     /*  Task object              */
int     result;   /*  completion result code   */
object  r;        /*  Task object              */
@end group
@end example
Advances to the next task in the parent @code{TaskList} with the given
result code.
@end deffn



@deffn {SetTag} gSetTag::Task
@findex gSetTag::Task
@sp 2
@example
@group
prev = gSetTag(task, tag);

object  task;    /*  Task object              */
object  tag;     /*  user data object         */
object  prev;    /*  previous tag or NULL     */
@end group
@end example
Attaches a user data object to the task.  Returns the previous tag.
@sp 1
See also:  @code{GetTag}, @code{AutoDisposeTag}
@end deffn



@deffn {TaskList} gTaskList::Task
@findex gTaskList::Task
@sp 2
@example
@group
tl = gTaskList(task);

object  task;   /*  Task object      */
object  tl;     /*  TaskList object   */
@end group
@end example
Returns the @code{TaskList} that this task belongs to.
@end deffn




@section HTTP Request
The @code{HttpRequest} class provides an HTTP client for making
HTTP and HTTPS requests.  It wraps the WinInet API and supports
custom headers, authentication, request bodies, and response handling.
The @code{WebService} class builds on @code{HttpRequest} to provide
a higher-level SOAP web service client.

@subsection HttpRequest Methods



@deffn {AddHeader} gAddHeader::HttpRequest
@findex gAddHeader
@sp 2
@example
@group
r = gAddHeader(req, name, value);

object  req;     /*  HttpRequest object  */
char    *name;   /*  header name         */
char    *value;  /*  header value        */
object  r;       /*  HttpRequest object  */
@end group
@end example
Adds a custom HTTP request header.
@end deffn



@deffn {GetErrorMessage} gGetErrorMessage::HttpRequest
@findex gGetErrorMessage::HttpRequest
@sp 2
@example
@group
s = gGetErrorMessage(req);

object  req;   /*  HttpRequest object       */
char    *s;    /*  error message or NULL    */
@end group
@end example
Returns the error message string if the request failed, or @code{NULL}
if no error occurred.
@end deffn



@deffn {GetResponse} gGetResponse::HttpRequest
@findex gGetResponse::HttpRequest
@sp 2
@example
@group
s = gGetResponse(req);

object  req;   /*  HttpRequest object    */
char    *s;    /*  response body string  */
@end group
@end example
Returns the response body as a string.  Only populated if no output
file was specified via @code{gSetOutputFileName}.
@sp 1
See also:  @code{SetOutputFileName}
@end deffn



@deffn {GetResponseHeader} gGetResponseHeader::HttpRequest
@findex gGetResponseHeader::HttpRequest
@sp 2
@example
@group
s = gGetResponseHeader(req, name);

object  req;    /*  HttpRequest object         */
char    *name;  /*  header name (case-insensitive)  */
char    *s;     /*  header value or NULL        */
@end group
@end example
Returns the value of a single response header by name.  The lookup is
case-insensitive.  Returns @code{NULL} if the header is not found.
@sp 1
See also:  @code{GetResponseHeaders}
@end deffn



@deffn {GetResponseHeaders} gGetResponseHeaders::HttpRequest
@findex gGetResponseHeaders
@sp 2
@example
@group
dict = gGetResponseHeaders(req);

object  req;    /*  HttpRequest object          */
object  dict;   /*  StringDictionary of headers */
@end group
@end example
Returns a @code{StringDictionary} containing all response headers.
Header names are lowercased.
@sp 1
See also:  @code{GetResponseHeader}
@end deffn



@deffn {NewHttpRequest} gNewHttpRequest::HttpRequest
@findex gNewHttpRequest
@sp 2
@example
@group
req = gNewHttpRequest(HttpRequest, url, verb);

object  req;    /*  new HttpRequest object  */
char    *url;   /*  request URL             */
char    *verb;  /*  HTTP method (GET, POST, etc.)  */
@end group
@end example
This is a class method.  Creates a new @code{HttpRequest} for the given
URL and HTTP method.
@end deffn



@deffn {SendRequest} gSendRequest::HttpRequest
@findex gSendRequest::HttpRequest
@sp 2
@example
@group
code = gSendRequest(req, flags);

object  req;    /*  HttpRequest object        */
DWORD   flags;  /*  connection flags          */
DWORD   code;   /*  0=success, or HTTP code   */
@end group
@end example
Sends the HTTP request with all configured headers and content.
Returns 0 on success (HTTP 2xx), or the HTTP response code on error.
@end deffn



@deffn {SetAuth} gSetAuth::HttpRequest
@findex gSetAuth
@sp 2
@example
@group
r = gSetAuth(req, user, pass);

object  req;    /*  HttpRequest object  */
char    *user;  /*  username            */
char    *pass;  /*  password            */
object  r;      /*  HttpRequest object  */
@end group
@end example
Sets HTTP Basic Authentication credentials for the request.
@end deffn



@deffn {SetContent} gSetContent::HttpRequest
@findex gSetContent
@sp 2
@example
@group
r = gSetContent(req, content);

object  req;      /*  HttpRequest object  */
char    *content; /*  request body        */
object  r;        /*  HttpRequest object  */
@end group
@end example
Sets the request body content as a string.  Replaces any previous
content.
@sp 1
See also:  @code{SetContentFromFile}
@end deffn



@deffn {SetContentFromFile} gSetContentFromFile::HttpRequest
@findex gSetContentFromFile
@sp 2
@example
@group
r = gSetContentFromFile(req, filename);

object  req;       /*  HttpRequest object         */
char    *filename; /*  file path                  */
object  r;         /*  HttpRequest object or NULL */
@end group
@end example
Loads the request body content from a file.  Returns @code{NULL} on
file error.
@sp 1
See also:  @code{SetContent}
@end deffn



@deffn {SetOutputFileName} gSetOutputFileName::HttpRequest
@findex gSetOutputFileName
@sp 2
@example
@group
r = gSetOutputFileName(req, name);

object  req;    /*  HttpRequest object  */
char    *name;  /*  file path or NULL   */
object  r;      /*  HttpRequest object  */
@end group
@end example
Sets an output file path to save the response body to instead of
keeping it in memory.  Pass @code{NULL} to disable file output.
@sp 1
See also:  @code{GetResponse}
@end deffn



@deffn {SetResponseTimeoutSeconds} gSetResponseTimeoutSeconds::HttpRequest
@findex gSetResponseTimeoutSeconds
@sp 2
@example
@group
gSetResponseTimeoutSeconds(req, seconds);

object  req;      /*  HttpRequest object  */
int     seconds;  /*  timeout in seconds  */
@end group
@end example
Sets the timeout for waiting on a server response.  The default is
60 seconds.
@end deffn



@deffn {SetSecure} gSetSecure::HttpRequest
@findex gSetSecure
@sp 2
@example
@group
r = gSetSecure(req, secure);

object  req;     /*  HttpRequest object  */
int     secure;  /*  non-zero for HTTPS  */
object  r;       /*  HttpRequest object  */
@end group
@end example
Enables or disables HTTPS for the request.  Pass a non-zero value
to use HTTPS (port 443 by default).
@end deffn



@subsection WebService Methods

The @code{WebService} class wraps @code{HttpRequest} to provide a
SOAP web service client.  It builds SOAP XML envelopes, sends them
via HTTP POST, and parses XML responses.



@deffn {AddArgument} gAddArgument::WebService
@findex gAddArgument
@sp 2
@example
@group
r = gAddArgument(ws, param, val);

object  ws;      /*  WebService object  */
char    *param;  /*  parameter name     */
char    *val;    /*  parameter value    */
object  r;       /*  WebService object  */
@end group
@end example
Adds a SOAP method parameter as a name-value pair.  Call multiple times
to build the parameter list.
@end deffn



@deffn {GetErrorMessage} gGetErrorMessage::WebService
@findex gGetErrorMessage::WebService
@sp 2
@example
@group
s = gGetErrorMessage(ws);

object  ws;   /*  WebService object        */
char    *s;   /*  error message or NULL    */
@end group
@end example
Returns the error message string if the request failed, or @code{NULL}
if no error occurred.
@end deffn



@deffn {GetResponse} gGetResponse::WebService
@findex gGetResponse::WebService
@sp 2
@example
@group
s = gGetResponse(ws);

object  ws;   /*  WebService object        */
char    *s;   /*  raw XML response string  */
@end group
@end example
Returns the raw XML response body from the SOAP service.
@sp 1
See also:  @code{GetXMLResponse}
@end deffn



@deffn {GetXMLResponse} gGetXMLResponse::WebService
@findex gGetXMLResponse
@sp 2
@example
@group
xml = gGetXMLResponse(ws);

object  ws;    /*  WebService object   */
object  xml;   /*  XMLNode object      */
@end group
@end example
Parses the response as XML and returns an @code{XMLNode} object for
structured access to the response elements.
@sp 1
See also:  @code{GetResponse}
@end deffn



@deffn {NewServiceRequest} gNewServiceRequest::WebService
@findex gNewServiceRequest
@sp 2
@example
@group
ws = gNewServiceRequest(WebService, url);

object  ws;    /*  new WebService object  */
char    *url;  /*  SOAP service URL       */
@end group
@end example
This is a class method.  Creates a new @code{WebService} request
targeting the given SOAP endpoint URL.
@end deffn



@deffn {SendRequest} gSendRequest::WebService
@findex gSendRequest::WebService
@sp 2
@example
@group
code = gSendRequest(ws, flags);

object  ws;     /*  WebService object         */
DWORD   flags;  /*  connection flags          */
DWORD   code;   /*  0=success, or HTTP code   */
@end group
@end example
Builds the SOAP XML envelope from the configured method and arguments,
sends it via HTTP POST, and retrieves the response.  Returns 0 on
success, or the HTTP response code on error.
@end deffn



@deffn {SetMethod} gSetMethod::WebService
@findex gSetMethod
@sp 2
@example
@group
r = gSetMethod(ws, method);

object  ws;       /*  WebService object  */
char    *method;  /*  SOAP method name   */
object  r;        /*  WebService object  */
@end group
@end example
Sets the SOAP method (operation) name to invoke on the remote service.
@end deffn



@deffn {SetNamespaceURL} gSetNamespaceURL::WebService
@findex gSetNamespaceURL
@sp 2
@example
@group
r = gSetNamespaceURL(ws, url);

object  ws;    /*  WebService object       */
char    *url;  /*  XML namespace URL       */
object  r;     /*  WebService object       */
@end group
@end example
Sets the XML namespace URL for the SOAP envelope.
@end deffn




@section Line Control
The @code{LineControl} class (a subclass of @code{Control}) implements a
control that draws a line on a dialog or window.  The line color and
style (solid, dotted, dashed) can be configured.

@subsection LineControl Methods



@deffn {New} New::LineControl
@findex vNew::LineControl
@sp 2
@example
@group
lc = vNew(LineControl, ctlID, name, dlg);

object  lc;       /*  new LineControl object  */
UINT    ctlID;    /*  control ID              */
char    *name;    /*  control name            */
object  dlg;      /*  dialog object           */
@end group
@end example
Creates a new @code{LineControl} associated with a control defined in
the dialog resource.  The line defaults to black and solid style.
@end deffn



@deffn {SetLineColor} gSetLineColor::LineControl
@findex gSetLineColor
@sp 2
@example
@group
r = gSetLineColor(lc, color);

object  lc;     /*  LineControl object  */
char    color;  /*  color index         */
object  r;      /*  LineControl object  */
@end group
@end example
Sets the line color by color index.
@sp 1
See also:  @code{SetLineStyle}
@end deffn



@deffn {SetLineStyle} gSetLineStyle::LineControl
@findex gSetLineStyle
@sp 2
@example
@group
r = gSetLineStyle(lc, style);

object  lc;     /*  LineControl object       */
char    style;  /*  line style (solid, dotted, dashed)  */
object  r;      /*  LineControl object       */
@end group
@end example
Sets the line style (solid, dotted, dashed, etc.).
@sp 1
See also:  @code{SetLineColor}
@end deffn




@section Rectangle Control
The @code{RectControl} class (a subclass of @code{Control}) implements a
control that draws a rectangle on a dialog or window.  It supports
configurable frame styles, fill patterns, background colors, and text
display with formatting.

@subsection RectControl Methods



@deffn {AllowEdge} gAllowEdge::RectControl
@findex gAllowEdge
@sp 2
@example
@group
prev = gAllowEdge(rc, val);

object  rc;     /*  RectControl object         */
int     val;    /*  1=allow edge, 0=no         */
int     prev;   /*  previous setting           */
@end group
@end example
Controls whether text extends to the edge of the rectangle.  Returns
the previous setting.
@end deffn



@deffn {Draw3D} gDraw3D::RectControl
@findex gDraw3D
@sp 2
@example
@group
prev = gDraw3D(rc, val);

object  rc;     /*  RectControl object         */
int     val;    /*  1=enable 3D, 0=disable     */
int     prev;   /*  previous setting           */
@end group
@end example
Enables or disables a 3D effect on the rectangle border.  Returns the
previous setting.
@end deffn



@deffn {New} New::RectControl
@findex vNew::RectControl
@sp 2
@example
@group
rc = vNew(RectControl, ctlID, name, dlg);

object  rc;       /*  new RectControl object  */
UINT    ctlID;    /*  control ID              */
char    *name;    /*  control name            */
object  dlg;      /*  dialog object           */
@end group
@end example
Creates a new @code{RectControl} associated with a control defined in
the dialog resource.  Defaults to a frame on all sides, no fill, and
black color.
@end deffn



@deffn {SetBackgroundColor} gSetBackgroundColor::RectControl
@findex gSetBackgroundColor
@sp 2
@example
@group
r = gSetBackgroundColor(rc, color);

object  rc;     /*  RectControl object  */
char    color;  /*  color index         */
object  r;      /*  RectControl object  */
@end group
@end example
Sets a solid background color for the rectangle interior.
@sp 1
See also:  @code{SetBackgroundRGBColor}, @code{SetFill}
@end deffn



@deffn {SetBackgroundColorRef} gSetBackgroundColorRef::RectControl
@findex gSetBackgroundColorRef
@sp 2
@example
@group
r = gSetBackgroundColorRef(rc, clr);

object    rc;    /*  RectControl object  */
COLORREF  clr;   /*  Windows COLORREF    */
object    r;     /*  RectControl object  */
@end group
@end example
Sets the background color using a Windows @code{COLORREF} value.
@sp 1
See also:  @code{SetBackgroundColor}, @code{SetBackgroundRGBColor}
@end deffn



@deffn {SetBackgroundRGBColor} gSetBackgroundRGBColor::RectControl
@findex gSetBackgroundRGBColor
@sp 2
@example
@group
r = gSetBackgroundRGBColor(rc, red, green, blue);

object  rc;     /*  RectControl object  */
int     red;    /*  red component       */
int     green;  /*  green component     */
int     blue;   /*  blue component      */
object  r;      /*  RectControl object  */
@end group
@end example
Sets the background color using RGB values.
@sp 1
See also:  @code{SetBackgroundColor}, @code{SetBackgroundColorRef}
@end deffn



@deffn {SetDT_Format} gSetDT_Format::RectControl
@findex gSetDT_Format
@sp 2
@example
@group
r = gSetDT_Format(rc, fmt);

object  rc;    /*  RectControl object        */
short   fmt;   /*  Windows DT_* format flags */
object  r;     /*  RectControl object        */
@end group
@end example
Sets text formatting flags using Windows @code{DT_*} constants
(e.g.@: @code{DT_LEFT}, @code{DT_CENTER}, @code{DT_WORDBREAK}).
@end deffn



@deffn {SetFill} gSetFill::RectControl
@findex gSetFill
@sp 2
@example
@group
r = gSetFill(rc, fill);

object  rc;    /*  RectControl object  */
char    fill;  /*  1=fill, 0=no fill  */
object  r;     /*  RectControl object  */
@end group
@end example
Enables or disables filling the rectangle interior.
@sp 1
See also:  @code{SetPattern}, @code{SetBackgroundColor}
@end deffn



@deffn {SetFrameColor} gSetFrameColor::RectControl
@findex gSetFrameColor
@sp 2
@example
@group
r = gSetFrameColor(rc, color);

object  rc;     /*  RectControl object  */
char    color;  /*  color index         */
object  r;      /*  RectControl object  */
@end group
@end example
Sets the frame (border) color.
@sp 1
See also:  @code{SetFrameStyle}, @code{SetFrameThickness}
@end deffn



@deffn {SetFrameStyle} gSetFrameStyle::RectControl
@findex gSetFrameStyle
@sp 2
@example
@group
r = gSetFrameStyle(rc, style);

object  rc;     /*  RectControl object                     */
char    style;  /*  frame style (which sides have borders) */
object  r;      /*  RectControl object                     */
@end group
@end example
Sets which sides of the rectangle have borders (left, right, top,
bottom, or combinations).
@sp 1
See also:  @code{SetFrameColor}, @code{SetFrameThickness}
@end deffn



@deffn {SetFrameThickness} gSetFrameThickness::RectControl
@findex gSetFrameThickness
@sp 2
@example
@group
r = gSetFrameThickness(rc, thickness);

object  rc;          /*  RectControl object        */
short   thickness;   /*  frame thickness in pixels  */
object  r;           /*  RectControl object         */
@end group
@end example
Sets the frame border thickness in pixels.
@sp 1
See also:  @code{SetFrameColor}, @code{SetFrameStyle}
@end deffn



@deffn {SetPattern} gSetPattern::RectControl
@findex gSetPattern
@sp 2
@example
@group
r = gSetPattern(rc, dotSize, stepSize, foreColor, backColor);

object  rc;         /*  RectControl object                 */
short   dotSize;    /*  dot size for pattern               */
short   stepSize;   /*  spacing between dots               */
char    foreColor;  /*  pattern foreground color index      */
char    backColor;  /*  pattern background color index      */
object  r;          /*  RectControl object                 */
@end group
@end example
Sets a dot pattern for filling the rectangle interior.  Must also call
@code{gSetFill} to enable filling.
@sp 1
See also:  @code{SetFill}, @code{SetBackgroundColor}
@end deffn



@deffn {SetStringValue} gSetStringValue::RectControl
@findex gSetStringValue::RectControl
@sp 2
@example
@group
r = gSetStringValue(rc, val);

object  rc;    /*  RectControl object  */
char    *val;  /*  display text        */
object  r;     /*  RectControl object  */
@end group
@end example
Sets the text displayed within the rectangle.
@sp 1
See also:  @code{StringValue}, @code{SetDT_Format}
@end deffn



@deffn {SetTextColor} gSetTextColor::RectControl
@findex gSetTextColor::RectControl
@sp 2
@example
@group
r = gSetTextColor(rc, color);

object  rc;     /*  RectControl object  */
char    color;  /*  color index         */
object  r;      /*  RectControl object  */
@end group
@end example
Sets the text color by color index.
@sp 1
See also:  @code{SetTextRGBColor}, @code{SetTextColorRef}
@end deffn



@deffn {SetTextColorRef} gSetTextColorRef::RectControl
@findex gSetTextColorRef
@sp 2
@example
@group
r = gSetTextColorRef(rc, clr);

object    rc;    /*  RectControl object  */
COLORREF  clr;   /*  Windows COLORREF    */
object    r;     /*  RectControl object  */
@end group
@end example
Sets the text color using a Windows @code{COLORREF} value.
@sp 1
See also:  @code{SetTextColor}, @code{SetTextRGBColor}
@end deffn



@deffn {SetTextRGBColor} gSetTextRGBColor::RectControl
@findex gSetTextRGBColor
@sp 2
@example
@group
r = gSetTextRGBColor(rc, red, green, blue);

object  rc;     /*  RectControl object  */
int     red;    /*  red component       */
int     green;  /*  green component     */
int     blue;   /*  blue component      */
object  r;      /*  RectControl object  */
@end group
@end example
Sets the text color using RGB values.
@sp 1
See also:  @code{SetTextColor}, @code{SetTextColorRef}
@end deffn



@deffn {StringValue} gStringValue::RectControl
@findex gStringValue::RectControl
@sp 2
@example
@group
s = gStringValue(rc);

object  rc;   /*  RectControl object   */
char    *s;   /*  display text string  */
@end group
@end example
Returns the current display text.
@sp 1
See also:  @code{SetStringValue}
@end deffn
